\chapter*{KẾT LUẬN}
\addcontentsline{toc}{chapter}{KẾT LUẬN}

\section*{Những kết quả đã đạt được}

Đồ án đã thiết kế, hiện thực hóa và đánh giá thực nghiệm một hệ thống truy hồi thông tin và hỏi đáp theo kiến trúc RAG kết hợp agent cho lĩnh vực pháp luật giao thông đường bộ Việt Nam, hoàn thành bốn mục tiêu trọng tâm:

\begin{enumerate}
    \item \textbf{Xây dựng luồng nạp và phân tích cấu trúc văn bản quy phạm pháp luật tất định:} Chuyển đổi 14 văn bản quy phạm pháp luật (ba luật hợp nhất, năm nghị định gồm Nghị định 168/2024/NĐ-CP và các nghị định sửa đổi, năm thông tư và quy chuẩn QCVN 41:2024/BGTVT) thành cây phân cấp Điều--Khoản--Điểm với 6.775 chunk và 10.570 cạnh viện dẫn nội bộ. Bộ phân tích hoạt động hoàn toàn dựa trên quy ước soạn thảo văn bản quy phạm pháp luật, loại trừ tự động các đoạn văn bản nằm trong dấu ngoặc kép của câu lệnh sửa đổi và kiểm chứng bất biến phủ dòng trên toàn bộ văn bản.
    \item \textbf{Xây dựng tầng truy hồi kết hợp có kiểm soát hiệu lực thời gian và ghi chú sửa đổi:} Kết hợp truy hồi vector ngữ nghĩa (\texttt{Qwen3-Embedding-0.6B} cắt còn 512 chiều theo Matryoshka) với chỉ mục ngược BM25 trên âm tiết và bigram không dấu thông qua hàm hợp nhất RRF ($k = 20$), bổ sung nhánh đối sánh nguyên văn xếp hạng theo vector truy vấn. Hệ thống lọc điều khoản theo ngày xảy ra vi phạm ngay trong SQL và tự động gắn ghi chú sửa đổi (\texttt{amended\_by}) được suy diễn từ chính các câu lệnh sửa đổi ở cấp Khoản, Điểm.
    \item \textbf{Thiết kế cơ chế điều phối agent đa bước qua giao thức MCP kèm hậu kiểm bám nguồn:} Cung cấp bốn công cụ truy xuất chỉ đọc có cấu trúc qua máy chủ MCP, kết hợp bước tái nạp nguyên văn điều khoản từ cơ sở dữ liệu và hàm kiểm tra quy tắc \texttt{check\_grounding} nhằm ngăn chặn triệt để hiện tượng mô hình tự sinh số hiệu điều luật hoặc mức tiền phạt không có trong nguồn.
    \item \textbf{Thiết lập khung đánh giá thực nghiệm hai tầng tách biệt:} Xây dựng bộ dữ liệu đánh giá truy hồi (tách \texttt{dev} và \texttt{holdout}) và bộ dữ liệu đánh giá câu trả lời cuối gồm 81 câu hỏi theo văn phong người dùng thực tế, thực hiện thực nghiệm ablation có kiểm định thống kê ghép cặp McNemar và đối sánh các mô hình embedding.
\end{enumerate}

\section*{Trả lời các câu hỏi nghiên cứu}

\textbf{CH1 -- Biểu diễn cấu trúc văn bản quy phạm pháp luật.} Biểu diễn cây phân cấp Điều--Khoản--Điểm dựa trên quy ước soạn thảo bảo toàn đầy đủ nội dung quy phạm và các quan hệ viện dẫn nội bộ mà không cần quy tắc riêng cho từng văn bản: mọi dòng văn bản chuẩn thuộc đúng một chunk, khung tiền phạt ở câu dẫn của Khoản cha được truyền xuống từng Điểm con trong văn bản kèm ngữ cảnh, và các cụm viện dẫn neo bởi từ khóa \texttt{``này''} được giải chính xác trên cây cú pháp. Cơ chế phân tích câu lệnh sửa đổi theo phạm vi kế thừa ở cấp Khoản và Điểm xác định chính xác 109 câu lệnh sửa đổi với 183 đích tác động trên 366 chunk, khắc phục triệt để tình trạng gắn dư 395/686 chunk của phương pháp chỉ đối sánh tiêu đề Điều.

\textbf{CH2 -- Phương pháp truy hồi điều khoản.} Trên bộ 715 câu đánh giá, hệ thống truy hồi đầy đủ đạt tỷ lệ tìm thấy đáp án 93,4\% trên tập \texttt{dev} và 97,6\% trên tập \texttt{holdout} ($\text{Hit@1}$ đạt 58,6\% trên \texttt{dev} và 64,2\% trên \texttt{holdout}; $\text{MRR}$ đạt 0,692 và 0,751), với độ trễ trung vị 501 ms và phân vị p95 là 690 ms. Thành phần tạo ra mức cải thiện lớn nhất là việc thay thế hàm \texttt{ts\_rank} của PostgreSQL bằng chỉ mục BM25 trên âm tiết và bigram không dấu, giúp $\text{Hit@1}$ của riêng nhánh từ khóa trên tập \texttt{dev} tăng từ 13,7\% lên 57,8\% ($p = 8{,}4\cdot10^{-77}$). Phép hợp nhất RRF giữa vector và BM25 vượt trội có ý nghĩa thống kê so với từng nhánh đơn lẻ, trong khi nhánh đối sánh nguyên văn xếp theo vector câu hỏi đóng vai trò lưới an toàn giúp khôi phục thêm 17 câu khó ($p = 1{,}5\cdot10^{-5}$).

\textbf{CH3 -- Điều phối agent đa bước và chất lượng câu trả lời cuối.} Trên bộ 81 câu hỏi theo văn phong người dùng thực tế, việc để agent tự chuyển đổi câu hỏi khẩu ngữ, không dấu, viết tắt sang truy vấn pháp lý và gọi công cụ qua MCP giúp hệ thống đạt tỷ lệ đúng toàn phần 97,2\% trên các câu trong phạm vi (69/71 câu; đạt 100\% ở cả ba nhóm khẩu ngữ, không dấu và trang trọng), không từ chối nhầm câu nào (0\%), tỷ lệ số liệu bám nguồn đạt 98,8\% và không có bất kỳ câu trả lời nào bịa mức tiền phạt, đồng thời từ chối đúng 80,0\% (8/10) số câu hỏi ngoài phạm vi. Phép đo ghép cặp trên 40 câu \texttt{holdout} cho thấy các bước gọi công cụ sau lần \texttt{search} đầu tiên không làm tăng tỷ lệ định vị đường dẫn (92,5\% so với 95,0\%), nhưng giữ vai trò quyết định trong việc đọc toàn văn Điều để lấy khung phạt, mở điều khoản sửa đổi và tự hiệu chỉnh số liệu bám nguồn.

\section*{Hạn chế và hướng phát triển}

Nghiên cứu hiện còn một số giới hạn: bộ đáp án chuẩn chưa qua bước thẩm định độc lập của chuyên gia pháp lý; kho dữ liệu mới bao phủ 14 văn bản và chưa bóc tách cấu trúc chi tiết bên trong các biểu mẫu hành chính; chế độ agent còn phụ thuộc vào mô hình ngôn ngữ thương mại với độ trễ trung vị 36 giây; và bộ suy diễn sửa đổi chưa giải các câu lệnh sửa đổi cấp Chương/Mục hoặc chuỗi văn bản sửa đổi bắc cầu.

Trong giai đoạn tiếp theo, hướng phát triển của đề tài tập trung vào bốn nội dung: (1) thẩm định bộ đáp án với người có chuyên môn pháp lý và chuẩn hóa tiêu chí đánh giá đối với nhóm câu hỏi ngoài phạm vi có liên đới pháp luật dân sự/hình sự; (2) thử nghiệm các mô hình xếp hạng lại (reranker) có cửa sổ ngữ cảnh dài và các mô hình ngôn ngữ mở kích thước nhỏ làm agent điều phối công cụ MCP nhằm xử lý 39 câu truy hồi còn trượt ở thứ hạng 11--40 đồng thời giảm độ trễ suy luận; (3) mở rộng bộ suy diễn sửa đổi cho các câu lệnh tác động tới cấp Chương/Mục và chuỗi văn bản sửa đổi bắc cầu; và (4) mở rộng kho văn bản quy phạm pháp luật kết hợp đánh giá thực tế với người dùng cuối.
